\section{Related Work}\label{sec:Related}
Static analysis for memory safety has advanced steadily, with recent work covering more complex scenarios. Multiple approaches have been developed for detecting use-after-free bugs: Teodorescu and Lucanu~\cite{static_analysis_use_after_free_cpp} developed frameworks for detecting use-after-free bugs in C++, Yan et al.~\cite{spatio_temporal_context_reduction} introduced spatio-temporal context reduction techniques for scalable UAF detection in multi-MLOC C code, and Caballero et al.~\cite{static_analysis_dangling_pointer} proposed early detection of dangling pointers. Recent comparative studies~\cite{hassler2025comparativestudyfuzzersstatic} evaluate both static analyzers and fuzzers for finding memory unsafety, showing complementarity across tools. These approaches largely focus on explicit memory management rather than conservative-GC settings where subordinate pointers can outlive visible roots.

Closest to our work is Ugawa and Fujimoto's study on detecting mistakes in registering local variables for precise garbage collection in an embedded JavaScript VM (eJSVM) using Coccinelle's control-flow--based pattern matching~\cite{UgawaFujimoto2019PRO}. While the setting is similar, in the case of \texttt{RB\_GC\_GUARD} we must track both derived-pointer extraction and derived-pointer usage, and we must reason about the absence of \texttt{VALUE}-level accesses after potential trigger points.

A closely related line of work in industrial language runtimes is static checking of GC rooting hazards. For example, SpiderMonkey (Mozilla's JavaScript engine) runs a ``rooting hazard analysis'' that reports stack-rooting hazards and related GC safety issues, and it is designed to be used both locally and in continuous integration\cite{spidermonkey_hazard_analysis}. While SpiderMonkey's setting differs (moving GC and explicit rooting APIs), the engineering motivation is similar: rooting/guarding mistakes are difficult to test reliably, and targeted static analysis can make these hazards reviewable at scale.

Empirical evaluations have shown CodeQL's potential for security bug detection, with studies revealing it can detect bug categories missed by other tools~\cite{charoenwet2024empirical}. CodeQL builds upon the foundational QL language~\cite{ql_ecoop} to provide high-level analysis capabilities, and has been used as the basis for domain-specific analyzers such as a resource leak checker for C\# code~\cite{resourceleak_csharp_codeql}. Youn et al.~\cite{codeql_multilingual} extended CodeQL for multilingual program analysis, tracking dataflows across language boundaries.

Recent work has also explored integrating large language models with static analysis. Li et al.\ propose LLift, a framework that uses LLMs to guide static analysis for use-before-initialization bugs in the Linux kernel~\cite{static_analysis_llm}. Paralegal~\cite{paralegal_static_analysis_privacy_bugs} demonstrates that carefully engineered static analyses can be made practical for privacy-bug detection at scale. These lines of work are complementary to GuardLint: they suggest that query-based analyses expressed in CodeQL or similar languages can serve as a foundation for richer, possibly LLM-augmented reasoning about subtle safety properties.
